\section{Guiding Applications}
\label{sec:1.2}

This textbook is primarily concerned with Monte Carlo methods for computing integrals of the form
\begin{equation}
  \calI_f[h] := \int_E h(x)\,f(x)\,dx = \Expect_{X \sim f}\!\left[h(X)\right],
  \label{eq:1.1}
\end{equation}
where $f$ is a given \textit{target distribution} supported on a set $E \subset \R^d$ and
$h : E \to \R$ is a given \textit{test function}. For instance, taking $h(x) = \mathbf{1}_A(x)$ to be
the indicator function of a subset $A \subset E$, we can then compute the probability
$\Prob_{X \sim f}(X \in A)$. In this section, we describe two application domains where computing
integrals of the form~\eqref{eq:1.1} is important: Bayesian statistics and statistical mechanics.
In both cases, the dimension $d$ can be exceedingly large and the target distribution $f$ is known
only up to a normalizing constant, rendering deterministic quadrature methods infeasible and posing
significant challenges for the design of Monte Carlo algorithms.

\paragraph{Bayesian Statistics}
Given a likelihood model $f(y|\theta)$ and a prior density $f(\theta)$, the posterior
density of the parameter $\theta$ given the data $y$ is
\begin{equation}
  f(\theta|y) = \frac{1}{c}\,f(y|\theta)\,f(\theta),
  \label{eq:1.2}
\end{equation}
where $c = \int f(y|\theta)\,f(\theta)\,d\theta$ is a normalizing constant ensuring that $f(\theta|y)$
integrates to one. Bayesian inference often requires computing expectations of the form
\[
  \calI_f[h] = \int h(\theta)\,f(\theta|y)\,d\theta,
\]
where the posterior $f(\theta|y)$ is the target distribution and $h$ is a test function. The most
important choice of test function is arguably $h(\theta) = \theta$, which gives the posterior mean
estimator of $\theta$. Other choices of test function enable computation of credible intervals
---regions of the parameter space with prescribed posterior probability--- as well as higher posterior
moments used for uncertainty quantification.

Before the advent of Monte Carlo methods, Bayesian inference was computationally feasible only in
low-dimensional settings or under the restrictive model assumption of conjugate priors: given a
likelihood, choosing a conjugate prior ensures that the posterior belongs to the same family of
distributions as the prior, whereby point estimators and credible intervals may be computed in
closed form. In modern applications, however, non-conjugate models are common and the parameter
$\theta$ can be high-dimensional; think, for instance, that $\theta$ may represent a vectorized
high-resolution image that we wish to denoise. Then, the normalizing constant $c$ in~\eqref{eq:1.2}
---known as the \textit{marginal likelihood} of the data--- is typically unknown and hard to compute,
as it is itself defined as an integral over a high-dimensional parameter space. This textbook therefore
emphasizes the importance of designing Monte Carlo methods that scale well to high dimension and that
are applicable when the normalizing constant of the target distribution is unknown.

The development of Monte Carlo methods goes hand in hand with the popularization of large-scale
Bayesian statistics. The Bayesian framework facilitates the formulation of hierarchical models and
enables sequential inference through the updating of beliefs as new data arrive. Gibbs samplers and
particle filters ---two classes of Monte Carlo methods studied in this textbook--- dovetail naturally
with hierarchical and sequential inference, respectively.

\paragraph{Statistical Mechanics}
Another important source of challenging sampling problems arises in statistical mechanics,
particularly in the simulation of molecular dynamics. Boltzmann postulated that the positions $q$
and momenta $p$ of the atoms in a molecular system of constant size, occupying a constant volume,
and in contact with a heat bath (at constant temperature), are distributed according to
\[
  f(q,p) = \frac{1}{c}\exp\!\left(-\beta\bigl(V(q) + K(p)\bigr)\right),
\]
where $c$ is a normalizing constant known as the \textit{partition function}, $\beta$ denotes the
inverse temperature, $V$ is a potential energy describing interactions between the particles in the
system, and $K$ represents the kinetic energy of the system. The potential $V$ is often a highly
irregular function with many local minima. This feature, along with other characteristics typical of
molecular systems, makes it difficult to compute averages with respect to $f(q,p)$. Nevertheless,
many quantities of scientific interest are defined as expectations under the Boltzmann distribution,
and computing integrals of the form
\[
  \calI_f[h] = \int h(q,p)\,f(q,p)\,dq\,dp.
\]
is therefore of central importance. As in Bayesian statistics, computing the normalizing constant $c$
of the Boltzmann distribution (also known as the Gibbs distribution) is challenging, which motivates
again the need to develop Monte Carlo methods for target densities that are only known up to a
normalizing constant.
